
\subsubsection{4.1 Implementation Details}

\begin{itemize}
\item \textbf{Model}: \texttt{cross-encoder/nli-deberta-v3-large} (304M parameters, frozen, float32)
\item \textbf{Hardware}: Single NVIDIA H100 NVL, CUDA\_VISIBLE\_DEVICES=0
\item \textbf{Batch size}: 32 (OOM fallback: 16); maximum sequence length: 512 tokens
\item \textbf{Training}: None; inference-only
\item \textbf{Total pairs evaluated}: 60,000 across three tasks
\item \textbf{Runtime}: Approximately 3.5 hours
\item \textbf{Unit tests}: 32/32 pass across all 5 test files (h-e1 experiment code)
\end{itemize}

\subsubsection{4.2 Baseline}

SelfCheckGPT-NLI \citep{Manakul2023SelfCheckGPT} on HaluEval: Dialogue AUROC = 0.48, QA AUROC = 0.530. Evaluated on base (non-instruction-tuned) Meta-Llama-3-8B, which produces near-uniform stochastic samples; this configuration likely underestimates SelfCheckGPT performance under intended deployment conditions.

\subsubsection{4.3 Mechanism Analysis Setup}

Statistical analysis for mechanism verification (h-m1) was conducted on the pre-computed score matrices from h-e1, using CPU-only computation (no GPU, no model re-inference). Tests: KL divergence from uniform (scipy.stats.entropy), Wilcoxon rank-sum (scipy.stats.ranksums). Total pairs analyzed: 60,000. Unit tests: 15/15 pass.
